\documentclass[conference]{IEEEtran}
\usepackage{physics}
\IEEEoverridecommandlockouts
\usepackage{blindtext}
\usepackage{amsmath}
\usepackage{gensymb}
\usepackage{graphicx}
\usepackage{xfrac}
\usepackage{float}
\graphicspath{ {./images/}  }

\begin{document}

\title{ECSE 331\\
    Laboratory No. 4\\
    Design of a BJT Amplifier in PSPICE\\
    \Large \today\\
}

\author
{
    \IEEEauthorblockN{Seyed Yazdan Zinati}
    \IEEEauthorblockA{260928077 \\
    \textit{seyed.zinati@mail.mcgill.ca}}

    \and

    \IEEEauthorblockN{James Wu}
    \IEEEauthorblockA{260929944 \\
    \textit{james.wu4@mail.mcgill.ca}}
}
\maketitle
\clearpage

\begin{abstract}

\end{abstract}

\section{Introduction}

\section{Analysis}

\section{Conclusion}
\clearpage

\section{Appendices}
\subsection{Small Signal Model (AC Analysis)}
\noindent From the small signal model we can deduce:
\begin{equation}
    R{in} = R_1 || R_2 || r_\pi
\end{equation}
\vspace{0.3cm}

\noindent$r_\pi = \sfrac{V_T}{I_B}$ and $V_T = 25\times 10^{-3}V$ therefore:
\begin{equation}
    I_B = \frac{V_T}{r_\pi} = \frac{25 \times 10^{-3}}{r_\pi}
\end{equation}
\vspace{0.3cm}

\noindent Write $V_{out}$ as a function of $i_b$ and $i_b$ as a function of $V_{in}$:
\begin{equation}
    V_{out} = -\beta i_b (R_C || R_L)
\end{equation}
\begin{equation}
    i_b = \sfrac{V_{in}}{r_{\pi}}
\end{equation}
\vspace{0.3cm}


\noindent From (3) and (4) we can write $V_{out}$ as follows:
\begin{equation}
    V_{out} = -\beta \frac{V_{in}}{r_\pi} (R_C || R_L)
\end{equation}
\vspace{0.3cm}

\noindent Rewriting Equation (5) we get the following:
%
\begin{equation}
    \frac{V_{out}}{V_{in}} = \abs{\frac{-\beta}{r_\pi} (R_C || R_L)} \Rightarrow \frac{V_{out}}{V_{in}} = \frac{\beta}{r_\pi} (R_C || R_L)
\end{equation}\\



\subsection{(Small Signal Model) DC Analysis}
\begin{equation}
    V_B - V_E = 0.7V
\end{equation}

\begin{equation}
    \frac{V_{CC} - V_B}{R_1} = I_B + \frac{V_B}{R_2} \Rightarrow V_B = \frac{-I_B R_1 R_2 - V_{CC} R_2}{R_1 + R_2}
\end{equation}
\vspace{0.5cm}

\noindent Choosing the following parameters:
$
\begin{cases}
    R_1 = R_2 = 10K\Omega\\
    R_{sig} = 1K\Omega\\
    \beta = 100
\end{cases}$ and From Equation (1) we find $r_\pi$:\\
\[R_{in} = R_1 || R_2 || r_\pi = 1K\Omega \Rightarrow r_\pi = 1250\Omega\]\\

\noindent From Equation (2) we find $I_B$:
\[I_B = \frac{25 \times10^{-3}}{1250} = 2 \times 10^{-5}A\]\\

\noindent Finding $R_C$ using (6) and knowing that $R_L = 10K\Omega$:\\
\[\frac{V_{out}}{V_{in}} = \frac{100(R_C || R_L)}{r_\pi} = 50 \Rightarrow \boxed{R_C = 666.67\Omega}\]\\

\noindent Using Equation (8) with $V_{CC} = 9V$, we find $V_B = 4.4V$ so from (7), we get $V_E = 3.7V$.

\[I_E = (\beta + 1)I_B \Rightarrow I_E = 2.02 \times 10^{-3}A\]
\[R_E = \frac{V_E}{I_E} \Rightarrow \boxed{R_E = 1831.68\Omega}\]

\vspace{2cm}

Resistor values to be used in our design:\\
\center
\begin{cases}
    R_1 = 10K\Omega\\
    R_2 = 10K\Omega\\
    R_{sig} = 1K\Omega\\
    RC = 666.67 \Omega\\
    R_E = 1831.68\Omega
\end{cases}
\end{document}
